\section{Introduction}
\label{sec:intro}

Distributed optimization creates parallel candidate updates every round. The
usual aggregation rule treats those candidates as exchangeable gradient
estimates and averages them. Hyperparameter optimization (HPO), in contrast,
usually treats candidates as separate jobs: run several trajectories, choose
one by validation performance, and throw away the others. These two practices
leave a gap. Distributed training has already paid for parallel exploration,
but standard aggregation does not use it as exploration; external HPO uses
exploration, but only by spending separate trajectories.

We ask whether HPO can be made an aggregation primitive. At each round, the
system deliberately creates a small menu of update candidates: different local
learning rates, outer optimizer controls, or other low-dimensional optimizer
arms. A cheap probe batch evaluates the candidate next states. The server then
selects or averages those candidates into a single shared model. We call this
family \emph{deliberately diverse aggregation} (\method). The central idea is
not merely to run HPO in parallel. It is to reinvest the evidence from parallel
arms into the model being trained, rather than selecting a separate winning
trajectory after the fact.

This distinction matters most under a fixed-compute comparison. If a cluster
is split into independent HPO trials, each parameter setting receives fewer
workers or fewer total update opportunities. In \method, all candidates are
scored during the same distributed round and the chosen or weighted update
advances one trajectory. Some arms can be formed directly from a common worker
update pool, which avoids splitting workers across arms altogether.

There is also a sharp null case. If all workers are identical, data are IID,
and arm labels are artificial, then the labels contain no persistent
hyperparameter information. Probe weights may still move due to realized
mini-batch noise, and this can sometimes filter stochastic updates, but it is
not HPO in the usual sense. The useful regime is deliberate
non-exchangeability: the system assigns meaningfully different optimizer arms,
the probe observes which candidate state is better for the present round, and
the aggregation rule turns that signal into a shared update.

Our contributions are:

\begin{enumerate}
\item We formulate deliberately diverse aggregation as a way to place
exploration and exploitation inside distributed optimization.
\item We give an exchangeability null showing why identical IID workers cannot
provide persistent HPO information from labels alone.
\item We describe worker-arm and server-control mechanisms that score
candidate next states with low-cost probes and then aggregate or select them.
\item We give fixed-compute cost accounting for probe overhead, emphasizing
when internal HPO is efficient rather than a disguised external search.
\item We report controlled local-SGD, tiny language-model, and multi-domain
synthetic audits. The main fixed-compute outer-control experiments beat
validation-selected external HPO, and the adaptive shift experiments show HPO
wins over external fixed-arm selection across language-model and multi-domain
settings.
\end{enumerate}

The paper is intentionally scoped. It is not a claim that validation probes
automatically improve every distributed run. We include negative and boundary
cases where exchangeability, naive grouping, or stale defaults prevent a clear
HPO win. The claim is narrower and, we think, more useful: when the distributed
round is deliberately diversified and the probe is cheap enough, aggregation
can become an efficient explore-exploit mechanism.
